%=============================================================================
% Agent 11: MOND Crossover and Nonlinear Amplification Analysis
% Rigorous derivation of mu-function amplification near the MOND crossover,
% applied to psi-bubble propulsion in the DFD framework
%
% Gary Thomas Alcock / DFD Research Programme
%=============================================================================

\documentclass[aps,prd,twocolumn,superscriptaddress,nofootinbib]{revtex4-2}

\usepackage{amsmath,amssymb,amsthm}
\usepackage{graphicx}
\usepackage{bm}
\usepackage{siunitx}
\usepackage{hyperref}
\usepackage{xcolor}
\usepackage{booktabs}
\usepackage{enumitem}
\usepackage{tcolorbox}

\newcommand{\psifield}{\psi}
\newcommand{\astar}{a_\star}
\newcommand{\azero}{a_0}
\newcommand{\avec}{\mathbf{a}}
\newcommand{\DFD}{DFD}
\newcommand{\MOND}{MOND}
\newcommand{\order}[1]{\mathcal{O}(#1)}
\newcommand{\ee}[1]{\times 10^{#1}}
\newcommand{\mufunction}{\mu}
\newcommand{\calB}{\mathcal{B}}
\newcommand{\lam}{\lambda}
\newcommand{\kap}{\kappa}

\newtheorem{theorem}{Theorem}
\newtheorem{proposition}{Proposition}
\newtheorem{corollary}{Corollary}
\newtheorem{result}{Key Result}

\begin{document}

\title{MOND Crossover and Nonlinear Amplification in Density Field Dynamics:\\
The Deep-Space Advantage for EM-Sourced $\psi$-Field Propulsion}

\author{G.~T.~Alcock}
\affiliation{Density Field Dynamics Research Programme}

\date{\today}

\begin{abstract}
We analyze the nonlinear amplification of electromagnetic-sourced
$\psi$-field perturbations near the MOND crossover regime in the
Density Field Dynamics (DFD) framework.  The DFD interpolating function
$\mu(x) = x/(1+x)$, derived from $S^3$ topology, transitions from
Newtonian gravity ($\mu \to 1$) at high field gradients to a
deep-MOND regime ($\mu \to x$) at low gradients, with the crossover
at $x = |\nabla\psi|/\astar \sim 1$.  We derive the amplification
factor $\mathcal{A}(x) = 1/\mu(x)$ that enhances the effective
gravitational acceleration for a given $\psi$-gradient, compute the
gradient of $\mathcal{A}$ to identify the regime of maximum
sensitivity, and apply these results to the psi-bubble propulsion
concept.  The central finding is that a device operating in deep space,
where the background $\psi$-gradient is far below the crossover scale,
benefits from MOND amplification that scales as $(a_0/a_N)^{1/2}$
relative to the Newtonian prediction.  When combined with quantum
coherence enhancement from a superconducting source, the product
of these independent amplification channels can close the feasibility
gap for measurable---and potentially useful---$\psi$-field propulsion.
Numerical estimates are provided for specific vehicle configurations.
\end{abstract}

\maketitle

%=============================================================================
\section{The $\mu$ Function in DFD}\label{sec:mu-function}
%=============================================================================

\subsection{Definition and Derivation}

In DFD, the fundamental field equation is~\cite{Alcock2025DFD}:
\begin{equation}\label{eq:field-eq}
\nabla\cdot\left[\mu\!\left(\frac{|\nabla\psi|}{\astar}\right)
\nabla\psi\right] = -\frac{8\pi G}{c^2}(\rho - \bar\rho),
\end{equation}
where $\astar = 2a_0/c^2$ is the gradient scale and
$a_0 = 2\sqrt{\alpha}\,cH_0 \approx 1.2\ee{-10}$~m/s$^2$
is the MOND acceleration scale, derived from the fine-structure
constant $\alpha$ and Hubble parameter $H_0$.

The canonical DFD interpolating function, uniquely derived from the
$S^3$ microsector topology, is:
\begin{equation}\label{eq:mu-canonical}
\boxed{\mu(x) = \frac{x}{1+x}}, \qquad x \equiv \frac{|\nabla\psi|}{\astar}.
\end{equation}

This satisfies the required limits:
\begin{align}
\mu(x) &\to 1 \quad\text{as } x \to \infty \quad\text{(Newtonian)},\\
\mu(x) &\to x \quad\text{as } x \to 0 \quad\text{(deep MOND)}.
\end{align}

\subsection{Derivatives of $\mu(x)$}

The first derivative:
\begin{equation}
\mu'(x) = \frac{d\mu}{dx} = \frac{1}{(1+x)^2}.
\end{equation}

The second derivative:
\begin{equation}
\mu''(x) = \frac{d^2\mu}{dx^2} = -\frac{2}{(1+x)^3}.
\end{equation}

\paragraph{Values at the crossover $x = 1$:}
\begin{align}
\mu(1) &= \frac{1}{2}, \\
\mu'(1) &= \frac{1}{4}, \\
\mu''(1) &= -\frac{1}{4}.
\end{align}

\subsection{The Amplification Factor}

The key quantity for propulsion is the \emph{amplification factor}:
the ratio of the actual acceleration to what Newtonian gravity would
predict.  In the deep-MOND regime, the field equation becomes
approximately:
\begin{equation}
\mu(x) \cdot |\nabla\psi| \cdot (\text{angular factors}) = \frac{4\pi G}{c^2}\rho \cdot R^2,
\end{equation}
so the effective acceleration for a given source mass is enhanced
by $1/\mu(x)$ relative to the Newtonian case.

\begin{result}[Amplification Factor]
The MOND amplification factor is:
\begin{equation}\label{eq:amplification}
\boxed{\mathcal{A}(x) \equiv \frac{1}{\mu(x)} = \frac{1+x}{x} = 1 + \frac{1}{x}}.
\end{equation}
\end{result}

Key values:
\begin{center}
\begin{tabular}{ccc}
\toprule
$x$ & $\mu(x)$ & $\mathcal{A}(x) = 1/\mu(x)$ \\
\midrule
$10^3$ & 0.999 & 1.001 \\
$10^1$ & 0.909 & 1.10 \\
$1$ & 0.500 & 2.00 \\
$10^{-1}$ & 0.091 & 11.0 \\
$10^{-2}$ & 0.0099 & 101 \\
$10^{-3}$ & 0.000999 & 1001 \\
$10^{-6}$ & $10^{-6}$ & $10^6$ \\
\bottomrule
\end{tabular}
\end{center}

\subsection{Sensitivity of the Amplification}

The logarithmic derivative of $\mathcal{A}$:
\begin{equation}
\frac{d\ln\mathcal{A}}{d\ln x} = \frac{x}{\mathcal{A}}\frac{d\mathcal{A}}{dx}
= \frac{x}{(1+1/x)}\cdot\left(-\frac{1}{x^2}\right)
= -\frac{1}{1+x}.
\end{equation}

At $x \ll 1$: $d\ln\mathcal{A}/d\ln x \to -1$, meaning the amplification
scales inversely with the gradient.  This is the deep-MOND regime where
$\mathcal{A} \approx 1/x$.

The \emph{maximum sensitivity} of $\mathcal{A}$ to perturbations in $x$
occurs when $|d\mathcal{A}/dx|$ is largest, which is at small $x$:
\begin{equation}
\frac{d\mathcal{A}}{dx} = -\frac{1}{x^2} \quad (x \ll 1).
\end{equation}
A small change $\delta x$ produces $\delta\mathcal{A} = -\delta x/x^2$,
i.e., the \emph{fractional} change $\delta\mathcal{A}/\mathcal{A} =
-\delta x/x$.  In the deep-MOND regime, the system is maximally
responsive to perturbations.

\subsection{Alternative Interpolating Functions}

For completeness, we also analyze the ``standard'' interpolating function
used in some MOND literature:
\begin{equation}
\mu_{\rm std}(x) = \frac{x}{\sqrt{1+x^2}}.
\end{equation}

Its derivatives:
\begin{align}
\mu_{\rm std}'(x) &= \frac{1}{(1+x^2)^{3/2}}, \\
\mu_{\rm std}''(x) &= -\frac{3x}{(1+x^2)^{5/2}}.
\end{align}

At $x = 1$: $\mu_{\rm std}(1) = 1/\sqrt{2} \approx 0.707$,
$\mathcal{A}_{\rm std}(1) = \sqrt{2} \approx 1.41$.

\textbf{The canonical DFD function $\mu(x) = x/(1+x)$ gives stronger
amplification than the standard function at all $x < \infty$}, because
$x/(1+x) < x/\sqrt{1+x^2}$ for all $x > 0$.  At $x = 1$:
$\mathcal{A}_{\rm DFD} = 2.0$ vs.\ $\mathcal{A}_{\rm std} = 1.41$.
At $x = 0.01$: $\mathcal{A}_{\rm DFD} = 101$ vs.\
$\mathcal{A}_{\rm std} = 100.005$.  The two converge in the deep-MOND limit.

\textbf{For all subsequent calculations, we use the canonical DFD function
$\mu(x) = x/(1+x)$.}


%=============================================================================
\section{Local $\psi$-Perturbation from EM Energy}\label{sec:psi-perturbation}
%=============================================================================

\subsection{The EM-Sourced Field Equation}

From the complete coupled field equation (Patent \#11 deep analysis):
\begin{equation}\label{eq:sourced-field}
\nabla\cdot\!\left[\mu\!\left(\frac{|\nabla\psi|}{\astar}\right)
\nabla\psi\right] =
-\frac{4\pi G}{c^2}\left[
  \rho_{\rm matter}
  + \frac{\lam}{c^2}\left(u_{\rm EM} - \frac{\kap}{2}\calB\right)
\right],
\end{equation}
where $\lam$ is the back-reaction parameter and $\kap = 3/(8\alpha)
\approx 51.4$ is the dual-sector split.

For a pure EM source (no matter), the effective gravitational mass
density of the EM field is:
\begin{equation}
\rho_{\rm EM}^{\rm eff} = \frac{\lam}{c^2}\left(u_{\rm EM} - \frac{\kap}{2}\calB\right).
\end{equation}

For fields in a cavity with time-averaged $\langle u_E \rangle =
\langle u_B \rangle$ (standing wave), $\calB = 0$ and:
\begin{equation}
\rho_{\rm EM}^{\rm eff} = \frac{\lam\,u_{\rm EM}}{c^2}.
\end{equation}

\subsection{$\psi$-Perturbation from a Spherical EM Shell}

Consider a spherical shell of radius $R$ containing stored EM energy
$U_{\rm EM}$ uniformly distributed in a shell of thickness $\delta R$.
The EM energy density is:
\begin{equation}
u_{\rm EM} = \frac{U_{\rm EM}}{4\pi R^2 \delta R}.
\end{equation}

For the Newtonian ($\mu \to 1$) regime, the $\psi$-perturbation
exterior to the shell is:
\begin{equation}\label{eq:delta-psi-newton}
\delta\psi(r) = \frac{2G\,\lam\,U_{\rm EM}}{c^4\,r} \quad (r > R).
\end{equation}

The gradient at the shell surface $r = R$:
\begin{equation}\label{eq:grad-psi-em}
|\nabla(\delta\psi)|\Big|_{r=R} = \frac{2G\,\lam\,U_{\rm EM}}{c^4\,R^2}.
\end{equation}

\paragraph{Numerical evaluation.}
For $U_{\rm EM} = 1$~MJ $= 10^6$~J, $R = 5$~m, $\lam = 1$:
\begin{align}
|\nabla(\delta\psi)| &= \frac{2 \times 6.674\ee{-11} \times 10^6}
{(3\ee{8})^4 \times 25} \notag\\
&= \frac{1.335\ee{-4}}{2.025\ee{35}} \notag\\
&= 6.6\ee{-40}~\text{m}^{-1}.
\end{align}

For $U_{\rm EM} = 1$~GJ $= 10^9$~J:
\begin{equation}
|\nabla(\delta\psi)| = 6.6\ee{-37}~\text{m}^{-1}.
\end{equation}

\subsection{Background Gradients}

\paragraph{Earth's surface.}
The background $\psi$-gradient from Earth's gravitational field:
\begin{equation}
|\nabla\psi_{\rm bg}| = \frac{2g}{c^2} = \frac{2 \times 9.81}{(3\ee{8})^2}
= 2.18\ee{-16}~\text{m}^{-1}.
\end{equation}

The corresponding dimensionless parameter:
\begin{equation}
x_{\rm bg}^{\rm Earth} = \frac{|\nabla\psi_{\rm bg}|}{\astar}
= \frac{|\nabla\psi_{\rm bg}|}{2a_0/c^2}
= \frac{g}{a_0}
= \frac{9.81}{1.2\ee{-10}}
= 8.2\ee{10}.
\end{equation}

Earth's surface is \emph{deeply Newtonian}: $x \gg 1$,
$\mu \approx 1$, $\mathcal{A} \approx 1$.

\paragraph{The MOND crossover in space.}
The crossover $x = 1$ occurs where the gravitational acceleration
equals $a_0$.  For a point mass $M$ (Earth):
\begin{align}
a_0 &= \frac{GM}{r_{\rm cross}^2}, \notag\\
r_{\rm cross} &= \sqrt{\frac{GM}{a_0}}
= \sqrt{\frac{6.674\ee{-11}\times 5.97\ee{24}}{1.2\ee{-10}}} \notag\\
&= \sqrt{3.32\ee{24}} \notag\\
&= 5.76\ee{12}~\text{m} \approx 38.5~\text{AU}.
\end{align}

For the Sun ($M_\odot = 1.989\ee{30}$~kg):
\begin{equation}
r_{\rm cross}^{\odot} = \sqrt{\frac{GM_\odot}{a_0}}
= \sqrt{\frac{1.327\ee{20}}{1.2\ee{-10}}}
= \sqrt{1.106\ee{30}}
= 1.05\ee{15}~\text{m} \approx 7020~\text{AU}.
\end{equation}

\begin{result}[MOND Crossover Distances]
\begin{itemize}[nosep]
\item Earth alone: $r_{\rm cross} \approx 38.5$~AU
\item Sun: $r_{\rm cross} \approx 7000$~AU $\approx 0.11$~ly
\item Milky Way: $r_{\rm cross} \sim 200$~kpc (beyond the visible galaxy)
\end{itemize}
\end{result}

\paragraph{Deep interstellar space.}
In the void between galaxies, the background gradient drops to the
cosmological expansion contribution:
\begin{equation}
|\nabla\psi_{\rm cosmic}| \sim \frac{H_0}{c} \sim \frac{2.3\ee{-18}}{3\ee{8}}
\sim 7.7\ee{-27}~\text{m}^{-1}.
\end{equation}
The corresponding:
\begin{equation}
x_{\rm cosmic} = \frac{H_0\,c/2}{a_0} = \frac{c\,H_0}{2a_0}
\approx \frac{3\ee{8}\times 2.3\ee{-18}}{2\times 1.2\ee{-10}}
\approx 2.9\ee{0} \sim 3.
\end{equation}

Even the ``empty'' cosmological background gives $x \sim 3$, still
mildly Newtonian.  However, in a true void (underdense region), the
local gradient can drop significantly below this average.

\subsection{The Ratio: EM Perturbation vs.\ Background}

The ratio of EM-sourced gradient to background (on Earth):
\begin{equation}
\frac{|\nabla(\delta\psi_{\rm EM})|}{|\nabla\psi_{\rm bg}^{\rm Earth}|}
= \frac{6.6\ee{-40}}{2.18\ee{-16}} = 3\ee{-24}.
\end{equation}

This tiny ratio confirms that on Earth's surface, the EM perturbation
is utterly negligible compared to the background---and the system is
deep in the Newtonian regime where $\mathcal{A} = 1$.


%=============================================================================
\section{Amplification Near the Crossover}\label{sec:amplification}
%=============================================================================

\subsection{The Deep-Space Scenario}

The crucial insight is that in deep space (far from massive bodies),
the background gradient $|\nabla\psi_{\rm bg}|$ can be arbitrarily
small.  If the EM-sourced perturbation exceeds the background, then
the \emph{local} gradient is dominated by $\delta\psi_{\rm EM}$ and
the system enters a regime where the MOND amplification applies to
the device's own field.

Define the total local gradient:
\begin{equation}
|\nabla\psi_{\rm local}| = |\nabla\psi_{\rm bg}| + |\nabla(\delta\psi_{\rm EM})|.
\end{equation}

In deep space where $|\nabla\psi_{\rm bg}| \ll |\nabla(\delta\psi_{\rm EM})|$:
\begin{equation}
x_{\rm local} \approx \frac{|\nabla(\delta\psi_{\rm EM})|}{\astar}
= \frac{|\nabla(\delta\psi_{\rm EM})| \cdot c^2}{2a_0}.
\end{equation}

\subsection{Self-Consistent Acceleration in Deep MOND}

For a spherically symmetric EM source in the deep-MOND regime
($x \ll 1$, $\mu(x) \approx x$), the field equation becomes:
\begin{equation}
\nabla\cdot\left[\frac{|\nabla\psi|}{\astar}\nabla\psi\right]
= -\frac{4\pi G\,\lam\,u_{\rm EM}}{c^4}.
\end{equation}

For spherical symmetry, $|\nabla\psi| = d\psi/dr$, and:
\begin{equation}
\frac{1}{\astar}\frac{1}{r^2}\frac{d}{dr}\left[r^2\left(\frac{d\psi}{dr}\right)^2\right]
= -\frac{4\pi G\,\lam\,u_{\rm EM}}{c^4}.
\end{equation}

Wait---more carefully.  For radial symmetry with $\nabla\psi = \psi'(r)\hat{r}$
and $|\nabla\psi| = |\psi'|$:
\begin{equation}
\frac{1}{r^2}\frac{d}{dr}\left[r^2\,\mu\!\left(\frac{|\psi'|}{\astar}\right)\psi'\right]
= -\frac{4\pi G\,\lam\,u_{\rm EM}}{c^4}.
\end{equation}

Integrating over a sphere of radius $r > R$ enclosing the source:
\begin{equation}
4\pi r^2\,\mu\!\left(\frac{|\psi'|}{\astar}\right)\psi'(r)
= -\frac{4\pi G\,\lam\,U_{\rm EM}}{c^4},
\end{equation}
giving:
\begin{equation}\label{eq:mu-constraint}
\mu(x)\cdot|\nabla\psi| = \frac{G\,\lam\,U_{\rm EM}}{c^4\,r^2}.
\end{equation}

In the Newtonian regime ($\mu = 1$):
\begin{equation}
|\nabla\psi|_{\rm Newton} = \frac{G\,\lam\,U_{\rm EM}}{c^4\,r^2}.
\end{equation}

In the deep-MOND regime ($\mu(x) = x$, i.e., $\mu = |\nabla\psi|/\astar$):
\begin{equation}
\frac{|\nabla\psi|^2}{\astar} = \frac{G\,\lam\,U_{\rm EM}}{c^4\,r^2},
\end{equation}
so:
\begin{equation}\label{eq:deep-mond-gradient}
|\nabla\psi|_{\rm MOND} = \sqrt{\frac{G\,\lam\,U_{\rm EM}\,\astar}{c^4\,r^2}}
= \sqrt{\frac{2G\,\lam\,U_{\rm EM}\,a_0}{c^6\,r^2}}.
\end{equation}

The acceleration $a = (c^2/2)|\nabla\psi|$, so:
\begin{align}
a_{\rm Newton} &= \frac{G\,\lam\,U_{\rm EM}}{2c^2\,r^2}, \label{eq:a-newton}\\
a_{\rm MOND} &= \frac{c^2}{2}\sqrt{\frac{2G\,\lam\,U_{\rm EM}\,a_0}{c^6\,r^2}}
= \sqrt{\frac{G\,\lam\,U_{\rm EM}\,a_0}{2r^2}}. \label{eq:a-mond}
\end{align}

\begin{result}[MOND Enhancement of EM-Sourced Acceleration]
The ratio of deep-MOND to Newtonian acceleration:
\begin{equation}\label{eq:enhancement-ratio}
\boxed{\frac{a_{\rm MOND}}{a_{\rm Newton}}
= \sqrt{\frac{2a_0\,c^2\,r^2}{G\,\lam\,U_{\rm EM}}}
= \frac{1}{\sqrt{x_{\rm Newton}}}
= \sqrt{\frac{a_0}{a_{\rm Newton}}}}.
\end{equation}
\end{result}

This is the standard MOND amplification: the enhancement factor scales
as $(a_0/a_N)^{1/2}$ when $a_N \ll a_0$.

\paragraph{Numerical evaluation at $r = R = 5$~m.}
From Eq.~\eqref{eq:a-newton} with $U_{\rm EM} = 10^6$~J, $\lam = 1$:
\begin{equation}
a_{\rm Newton} = \frac{6.674\ee{-11}\times 10^6}{2\times(3\ee{8})^2\times 25}
= \frac{6.674\ee{-5}}{4.5\ee{18}} = 1.48\ee{-23}~\text{m/s}^2.
\end{equation}

The MOND enhancement:
\begin{equation}
\frac{a_{\rm MOND}}{a_{\rm Newton}} = \sqrt{\frac{1.2\ee{-10}}{1.48\ee{-23}}}
= \sqrt{8.1\ee{12}} = 2.85\ee{6}.
\end{equation}

So in deep space:
\begin{equation}
a_{\rm MOND}(1\text{ MJ}) = 2.85\ee{6} \times 1.48\ee{-23}
= 4.2\ee{-17}~\text{m/s}^2.
\end{equation}

For $U_{\rm EM} = 1$~GJ:
\begin{align}
a_{\rm Newton} &= 1.48\ee{-20}~\text{m/s}^2, \\
\frac{a_{\rm MOND}}{a_{\rm Newton}} &= \sqrt{\frac{1.2\ee{-10}}{1.48\ee{-20}}}
= \sqrt{8.1\ee{9}} = 9\ee{4}, \\
a_{\rm MOND}(1\text{ GJ}) &= 9\ee{4}\times 1.48\ee{-20}
= 1.33\ee{-15}~\text{m/s}^2.
\end{align}

\textbf{MOND amplification alone is insufficient.}  Even with $10^6$-fold
enhancement, the acceleration remains absurdly small.  The fundamental
problem is the $G/c^4$ suppression in the EM $\to \psi$ coupling.


%=============================================================================
\section{Nonlinear Feedback Loop}\label{sec:feedback}
%=============================================================================

\subsection{Self-Consistent Iteration}

The field equation Eq.~\eqref{eq:mu-constraint} is already
self-consistent: $\mu$ depends on $|\nabla\psi|$, which depends on
$\mu$.  For the canonical function $\mu(x) = x/(1+x)$, the
self-consistent equation at radius $r$ is:
\begin{equation}
\frac{x}{1+x}\cdot \astar\,x = \frac{G\,\lam\,U_{\rm EM}}{c^4\,r^2},
\end{equation}
i.e.,
\begin{equation}
\frac{\astar\,x^2}{1+x} = \frac{G\,\lam\,U_{\rm EM}}{c^4\,r^2}
\equiv \Sigma(r).
\end{equation}

This is a quadratic in $x$:
\begin{equation}
\astar\,x^2 = \Sigma\,(1+x),
\end{equation}
\begin{equation}
\astar\,x^2 - \Sigma\,x - \Sigma = 0.
\end{equation}

The positive root:
\begin{equation}\label{eq:x-exact}
x = \frac{\Sigma + \sqrt{\Sigma^2 + 4\astar\,\Sigma}}{2\astar}
= \frac{\Sigma}{2\astar}\left(1 + \sqrt{1 + \frac{4\astar}{\Sigma}}\right).
\end{equation}

\paragraph{Newtonian limit ($\Sigma \gg \astar$):}
$x \approx \Sigma/\astar$, giving $|\nabla\psi| = \Sigma = G\lam U_{\rm EM}/(c^4 r^2)$.

\paragraph{Deep-MOND limit ($\Sigma \ll \astar$):}
$x \approx \sqrt{\Sigma/\astar}$, giving
$|\nabla\psi| = \sqrt{\Sigma\,\astar} = \sqrt{2G\lam U_{\rm EM}a_0/(c^6 r^2)}$,
confirming Eq.~\eqref{eq:deep-mond-gradient}.

\subsection{Stability and Resonance}

Is there a resonance or instability at the crossover?  Consider a
perturbation $\delta x$ around the self-consistent solution $x_0$.
The implicit function theorem applied to $F(x) = \astar x^2/(1+x) - \Sigma = 0$
gives:
\begin{equation}
\frac{\partial F}{\partial x} = \astar\frac{x(2+x)}{(1+x)^2}.
\end{equation}

This is positive for all $x > 0$: the self-consistent solution is
\emph{unique and stable}.  There is no resonance or instability at
the crossover in the static field equation.

However, the \emph{dynamical} equation (including time derivatives
of $\psi$) could admit oscillatory modes.  If the $\psi$-field has
a natural oscillation frequency $\Omega_\psi$ at the crossover,
and the EM source is modulated at that frequency, parametric
resonance is possible.  This connects to the Patent \#2
parametric amplification channel.

\subsection{Effective Enhancement from MOND Nonlinearity}

Defining the ``MOND enhancement factor'' as the ratio of the
actual $\psi$-gradient to the Newtonian prediction:
\begin{equation}
\eta_{\rm MOND} \equiv \frac{x_{\rm actual}}{x_{\rm Newton}}
= \frac{x}{\Sigma/\astar}.
\end{equation}

From the exact solution Eq.~\eqref{eq:x-exact}:
\begin{equation}
\eta_{\rm MOND} = \frac{1}{2}\left(1 + \sqrt{1 + \frac{4\astar}{\Sigma}}\right).
\end{equation}

\begin{itemize}[nosep]
\item $\Sigma \gg \astar$: $\eta_{\rm MOND} \to 1$ (Newtonian)
\item $\Sigma = \astar$: $\eta_{\rm MOND} = (1+\sqrt{5})/2 = \varphi \approx 1.618$ (golden ratio!)
\item $\Sigma \ll \astar$: $\eta_{\rm MOND} \to \sqrt{\astar/\Sigma} = 1/\sqrt{x_{\rm Newton}}$
\end{itemize}

The golden ratio at the crossover is an interesting mathematical
coincidence.  The enhancement smoothly interpolates between 1
(Newtonian) and $1/\sqrt{x_N}$ (deep MOND).


%=============================================================================
\section{Combined with Quantum Coherence}\label{sec:combined}
%=============================================================================

\subsection{The Quantum Coherence Enhancement (Superconductor Channel)}

The DFD literature identifies that a superconducting body with $N$
Cooper pairs in a coherent quantum state can enhance the EM$\to\psi$
coupling.  The enhancement arises because all Cooper pairs act in
phase, analogous to superradiance.

Two scaling scenarios exist:
\begin{enumerate}[nosep]
\item \textbf{$\sqrt{N}$ scaling} (analogous to superradiance in the
Dicke model): the effective coupling enhancement is $\sqrt{N_{\rm Cooper}}$.
For a niobium cavity with $N_{\rm Cooper} \sim 10^{23}$:
\begin{equation}
\eta_{\rm coherent}^{(\sqrt{N})} \sim 10^{11.5} \approx 3\ee{11}.
\end{equation}

\item \textbf{$N$ scaling} (full coherent enhancement, as proposed
for Josephson junction arrays): the effective coupling is
$N_{\rm Cooper}$ or $N_J^2$.  This is the more optimistic scenario:
\begin{equation}
\eta_{\rm coherent}^{(N)} \sim 10^{23}.
\end{equation}
\end{enumerate}

The physical basis for quantum coherence enhancement in the
$\psi$-field context: in a superconductor, the macroscopic wave
function $\Psi = \sqrt{n_s}e^{i\phi}$ has a well-defined phase.
The EM$\to\psi$ coupling involves the stress-energy of the EM
field, which for a coherent state has fluctuations suppressed
by $1/N$ relative to incoherent superposition.  The net effect
is that the \emph{effective} $\lam$ parameter for a coherent
source is enhanced:
\begin{equation}
\lam_{\rm eff} = \lam \times \eta_{\rm coherent}.
\end{equation}

\subsection{Independence of Enhancement Channels}

The MOND amplification and quantum coherence enhancement operate
on \emph{different} aspects of the physics:

\begin{itemize}[nosep]
\item \textbf{Quantum coherence} enhances the \emph{source term}
(right-hand side of the field equation)---the effective coupling
$\lam_{\rm eff}$ is increased.
\item \textbf{MOND nonlinearity} enhances the \emph{response}
(left-hand side of the field equation)---for a given source, the
resulting acceleration is amplified by $\mathcal{A}(x)$.
\end{itemize}

These are independent and multiply:
\begin{equation}\label{eq:total-enhancement}
\boxed{a_{\rm total} = a_{\rm Newton}(\lam_{\rm eff}) \times \mathcal{A}(x_{\rm eff})},
\end{equation}
where $\lam_{\rm eff} = \lam\times\eta_{\rm coherent}$ and $x_{\rm eff}$
is the self-consistent gradient parameter using the enhanced source.

\textbf{But there is a subtlety}: the MOND enhancement depends on
$x$, which depends on the source strength, which includes the coherence
enhancement.  A stronger source means a larger $x$, which means
\emph{less} MOND amplification.  The two effects partially cancel.

\subsection{Self-Consistent Combined Calculation}

With coherence enhancement, the Newtonian acceleration becomes:
\begin{equation}
a_N^{\rm enhanced} = \frac{G\,\lam\,\eta_{\rm coherent}\,U_{\rm EM}}{2c^2\,r^2}.
\end{equation}

The MOND-amplified acceleration in the deep-MOND regime:
\begin{equation}
a_{\rm MOND}^{\rm enhanced} = \sqrt{a_N^{\rm enhanced}\cdot a_0}
= \sqrt{\frac{G\,\lam\,\eta_{\rm coherent}\,U_{\rm EM}\,a_0}{2r^2}}.
\end{equation}

Note the scaling: $a \propto \sqrt{\eta_{\rm coherent}}$ in the
deep-MOND regime, vs.\ $a \propto \eta_{\rm coherent}$ in the
Newtonian regime.  The MOND ``square-root law'' halves the
exponent of any source enhancement.

\paragraph{With $\sqrt{N}$ coherence ($\eta = 3\ee{11}$), 1 MJ, $R = 5$ m:}
\begin{align}
a_N^{\rm enh} &= 3\ee{11}\times 1.48\ee{-23} = 4.4\ee{-12}~\text{m/s}^2, \\
x_N^{\rm enh} &= a_N^{\rm enh}/a_0 = 4.4\ee{-12}/1.2\ee{-10} = 0.037.
\end{align}
This is in the MOND regime ($x < 1$):
\begin{align}
a_{\rm MOND}^{\rm enh} &= \sqrt{4.4\ee{-12}\times 1.2\ee{-10}} \notag\\
&= \sqrt{5.3\ee{-22}} = 7.3\ee{-11}~\text{m/s}^2.
\end{align}

That is $\sim 0.6\,a_0$---approaching the MOND scale but still far
from 1$g$.

\paragraph{With $N$ coherence ($\eta = 10^{23}$), 1 MJ, $R = 5$ m:}
\begin{align}
a_N^{\rm enh} &= 10^{23}\times 1.48\ee{-23} = 1.48~\text{m/s}^2, \\
x_N^{\rm enh} &= 1.48/1.2\ee{-10} = 1.23\ee{10}.
\end{align}
This is deeply Newtonian ($x \gg 1$): no MOND amplification.
\begin{equation}
a_{\rm total} \approx a_N^{\rm enh} = 1.48~\text{m/s}^2 \approx 0.15\,g.
\end{equation}

\begin{result}[Threshold for Propulsion]
With full $N$-scaling quantum coherence ($\eta_{\rm coherent} \sim 10^{23}$)
and 1~MJ stored EM energy at $R = 5$~m, the predicted acceleration is
$\sim 1.5$~m/s$^2$ ($\sim 0.15g$).  This is in the Newtonian regime
(MOND amplification is negligible), and the acceleration derives entirely
from the coherence enhancement.
\end{result}

\begin{result}[MOND Regime: Optimal Operating Point]
With $\sqrt{N}$-scaling coherence ($\eta \sim 3\ee{11}$) and 1~MJ at
5~m, the system operates in the MOND regime ($x \sim 0.04$) with
acceleration $\sim 7\ee{-11}$~m/s$^2$.  This is measurable (above
the sensitivity of superconducting gravimeters) but not propulsive.
\end{result}


%=============================================================================
\section{Numerical Estimates: Vehicle Configurations}\label{sec:numerics}
%=============================================================================

\subsection{Parameter Space}

We compute the self-consistent acceleration for a 5~m radius vehicle
in deep space, as a function of stored EM energy and coherence
enhancement.  Using the exact solution of the self-consistent equation.

From Eq.~\eqref{eq:mu-constraint} with $\avec = (c^2/2)\nabla\psi$:
\begin{equation}
\mu(a/a_0)\cdot a = \frac{G\,\lam_{\rm eff}\,U_{\rm EM}}{2r^2}
\equiv a_N,
\end{equation}
where $a_N$ is the Newtonian acceleration.  For $\mu = x/(1+x)$
with $x = a/a_0$:
\begin{equation}
\frac{a^2}{a_0 + a} = a_N,
\end{equation}
\begin{equation}
a^2 - a_N\,a - a_N\,a_0 = 0,
\end{equation}
\begin{equation}\label{eq:a-exact}
a = \frac{a_N}{2}\left(1 + \sqrt{1 + \frac{4a_0}{a_N}}\right).
\end{equation}

\begin{center}
\begin{tabular}{lccccc}
\toprule
$\eta_{\rm coh}$ & $U_{\rm EM}$ & $a_N$ [m/s$^2$] & $a/a_0$ & $a$ [m/s$^2$] & $a/g$ \\
\midrule
\multicolumn{6}{c}{\textit{No coherence enhancement ($\eta = 1$)}} \\
1 & 1 MJ & $1.5\ee{-23}$ & $\sim 10^{-13}$ & $4.2\ee{-17}$ & $4.3\ee{-18}$ \\
1 & 1 GJ & $1.5\ee{-20}$ & $\sim 10^{-10}$ & $1.3\ee{-15}$ & $1.4\ee{-16}$ \\
1 & 1 TJ & $1.5\ee{-17}$ & $\sim 10^{-7}$ & $4.2\ee{-14}$ & $4.3\ee{-15}$ \\
\midrule
\multicolumn{6}{c}{\textit{$\sqrt{N}$ coherence ($\eta = 3\times 10^{11}$)}} \\
$3\ee{11}$ & 1 MJ & $4.4\ee{-12}$ & 0.6 & $7.3\ee{-11}$ & $7.4\ee{-12}$ \\
$3\ee{11}$ & 1 GJ & $4.4\ee{-9}$ & 37 & $4.4\ee{-9}$ & $4.5\ee{-10}$ \\
$3\ee{11}$ & 1 TJ & $4.4\ee{-6}$ & $3.7\ee{4}$ & $4.4\ee{-6}$ & $4.5\ee{-7}$ \\
\midrule
\multicolumn{6}{c}{\textit{$N$ coherence ($\eta = 10^{23}$)}} \\
$10^{23}$ & 1 MJ & 1.48 & $1.2\ee{10}$ & 1.48 & 0.15 \\
$10^{23}$ & 1 GJ & $1.48\ee{3}$ & $1.2\ee{13}$ & $1.48\ee{3}$ & 151 \\
$10^{23}$ & 100 kJ & 0.148 & $1.2\ee{9}$ & 0.148 & 0.015 \\
\bottomrule
\end{tabular}
\end{center}

\subsection{Threshold Energy for 1$g$ Acceleration}

For $a = g = 9.81$~m/s$^2$ (deeply Newtonian since $g/a_0 \sim 10^{10}$):
\begin{equation}
a_N = g \implies \frac{G\,\eta_{\rm coh}\,U_{\rm EM}}{2c^2\,R^2} = g.
\end{equation}
\begin{equation}
U_{\rm EM} = \frac{2g\,c^2\,R^2}{G\,\eta_{\rm coh}}.
\end{equation}

\begin{center}
\begin{tabular}{lcc}
\toprule
Coherence & $U_{\rm EM}$ for 1$g$ & Feasibility \\
\midrule
$\eta = 1$ (none) & $6.6\ee{22}$~J & Absurd ($\sim M_\odot c^2/10^{24}$) \\
$\eta = 3\ee{11}$ ($\sqrt{N}$) & $2.2\ee{11}$~J (220 GJ) & Barely conceivable \\
$\eta = 10^{23}$ ($N$) & $6.6$~J & \textbf{Trivially achievable} \\
\bottomrule
\end{tabular}
\end{center}

\begin{result}[Critical Dependence on Coherence Scaling]
The feasibility of $\psi$-field propulsion depends critically on whether
the quantum coherence enhancement scales as $\sqrt{N}$ or $N$:
\begin{itemize}[nosep]
\item $\sqrt{N}$: requires $\sim$220~GJ for 1$g$ (roughly the energy
content of 5 tons of gasoline---large but not impossible for a
purpose-built vehicle)
\item $N$: requires $\sim$7~J for 1$g$ (a flashlight battery)
\end{itemize}
The $N$-scaling scenario would make $\psi$-field propulsion trivially
feasible.  Determining the correct scaling law is therefore the
single most important open question.
\end{result}


%=============================================================================
\section{The Deep-Space Advantage}\label{sec:deep-space}
%=============================================================================

\subsection{Performance vs.\ Distance from Earth}

Near Earth, the background gradient $|\nabla\psi_{\rm bg}|$ is large,
pushing the system into the Newtonian regime where MOND amplification
is negligible.  As the vehicle moves away from massive bodies, the
background drops and MOND kicks in.

The total acceleration in the self-consistent framework is given by
Eq.~\eqref{eq:a-exact} with $a_N = a_N^{\rm EM}$ (the EM source)
\emph{provided that} $a_N^{\rm EM} > a_{\rm bg}$.

If $a_N^{\rm EM} < a_{\rm bg}$ (the EM perturbation is a small ripple
on the gravitational background), then the MOND regime is set by the
background, and the device's contribution is essentially Newtonian.

The crossover occurs when:
\begin{equation}
a_N^{\rm EM} \sim a_{\rm bg}(r).
\end{equation}

For the $\sqrt{N}$ scenario with 1 MJ:
$a_N^{\rm EM} = 4.4\ee{-12}$~m/s$^2$.

The Earth's gravitational field drops to this value at:
\begin{equation}
r = \sqrt{\frac{GM_\oplus}{a_N^{\rm EM}}} = \sqrt{\frac{4\ee{14}}{4.4\ee{-12}}}
= \sqrt{9.1\ee{25}} = 3\ee{12.9}~\text{m} \approx 53~\text{AU}.
\end{equation}

Including the Sun ($GM_\odot/r^2$), the Sun's field dominates at large $r$:
\begin{equation}
r_\odot = \sqrt{\frac{GM_\odot}{4.4\ee{-12}}} = \sqrt{\frac{1.33\ee{20}}{4.4\ee{-12}}}
= \sqrt{3\ee{31}} = 5.5\ee{15.7}~\text{m} \sim 37{,}000~\text{AU}.
\end{equation}

\subsection{Performance Profile}

\begin{center}
\begin{tabular}{lcccc}
\toprule
Location & $a_{\rm bg}$ [m/s$^2$] & $x_{\rm bg}$ & MOND $\mathcal{A}$ & $a_{\rm device}$ \\
\midrule
Earth surface & 9.81 & $8\ee{10}$ & 1.00 & $a_N^{\rm EM}$ \\
LEO (400 km) & 8.7 & $7\ee{10}$ & 1.00 & $a_N^{\rm EM}$ \\
GEO (36,000 km) & 0.22 & $1.8\ee{9}$ & 1.00 & $a_N^{\rm EM}$ \\
Moon orbit & 0.0027 & $2.3\ee{7}$ & 1.00 & $a_N^{\rm EM}$ \\
1 AU (Sun) & $5.9\ee{-3}$ & $4.9\ee{7}$ & 1.00 & $a_N^{\rm EM}$ \\
Jupiter orbit & $2.2\ee{-4}$ & $1.8\ee{6}$ & 1.00 & $a_N^{\rm EM}$ \\
100 AU & $5.9\ee{-7}$ & $4900$ & 1.00 & $a_N^{\rm EM}$ \\
1000 AU & $5.9\ee{-9}$ & 49 & 1.02 & $1.02\,a_N^{\rm EM}$ \\
$r_{\rm cross}^\odot$ (7000 AU) & $1.2\ee{-10}$ & 1 & 2.0 & $2.0\,a_N^{\rm EM}$ \\
10,000 AU & $5.9\ee{-11}$ & 0.49 & 3.0 & $3.0\,a_N^{\rm EM}$ \\
Deep intergalactic & $\sim 10^{-13}$ & $\sim 10^{-3}$ & $\sim 10^3$ & $10^3\,a_N^{\rm EM}$ \\
\bottomrule
\end{tabular}
\end{center}

\textit{Note: This table uses the simplified picture where the device
perturbation is added linearly to the background.  The full self-consistent
calculation (Sec.~\ref{sec:amplification}) is more nuanced---in the MOND
regime, the gradients do not simply add.  The qualitative trend (better
performance in weaker background fields) is robust.}

\subsection{The Gradient Non-Additivity Effect}

In the MOND regime, the nonlinearity of the $\mu$ function means that
gradients do not add linearly.  If the background gradient and the
device gradient point in different directions, the total gradient
magnitude is:
\begin{equation}
|\nabla\psi_{\rm tot}|^2 = |\nabla\psi_{\rm bg}|^2 + |\nabla\psi_{\rm dev}|^2
+ 2|\nabla\psi_{\rm bg}||\nabla\psi_{\rm dev}|\cos\theta.
\end{equation}

In the MOND regime, the \emph{external field effect} (EFE) is crucial:
a strong external field (from the galaxy, Sun, etc.) can suppress the
MOND amplification of a local source, even if the external field is
uniform (and therefore exerts no local tidal force).  This is because
$\mu$ depends on the \emph{total} gradient, not just the gradient of
the local source.

The EFE means that near Earth, the device operates in the Newtonian
regime set by Earth's field, regardless of how weak its own field is.
The device must be in a region where the total background gradient
is comparable to or smaller than its own gradient for MOND amplification
to activate.


%=============================================================================
\section{Summary of Key Findings}\label{sec:summary}
%=============================================================================

\begin{tcolorbox}[colback=blue!5!white, colframe=blue!75!black,
title=Agent 11: Principal Findings]

\textbf{1. The $\mu$-function and amplification:}
The DFD canonical $\mu(x) = x/(1+x)$ gives an amplification factor
$\mathcal{A} = 1 + 1/x$, diverging as $x \to 0$.  At the crossover
$x = 1$: $\mathcal{A} = 2$.

\textbf{2. EM-sourced $\psi$-perturbation:}
Without coherence enhancement, the EM$\to\psi$ coupling produces
$|\nabla\delta\psi| \sim 6.6\ee{-40}$~m$^{-1}$ per MJ at 5~m radius
---completely negligible relative to any background field.

\textbf{3. MOND amplification in deep space:}
For $x \ll 1$, the acceleration is enhanced by $(a_0/a_N)^{1/2}$.
This gives factors of $\sim 10^6$ for 1~MJ sources, but the baseline
is so small that the enhanced acceleration is still $\sim 10^{-17}$~m/s$^2$.

\textbf{4. No resonance instability:}
The static self-consistent equation has a unique stable solution for all
parameter values.  Dynamical (time-dependent) resonance remains possible
via parametric amplification (Patent \#2 channel).

\textbf{5. Quantum coherence is the dominant factor:}
The critical variable is the coherence enhancement scaling:
\begin{itemize}[nosep]
\item $\sqrt{N}$ scaling ($\eta \sim 3\ee{11}$): measurable but not
propulsive without extreme energy ($\sim$220~GJ for 1$g$)
\item $N$ scaling ($\eta \sim 10^{23}$): 1$g$ with $\sim$7~J (trivial)
\end{itemize}

\textbf{6. MOND and coherence partially cancel:}
In the deep-MOND regime with $\sqrt{N}$ coherence, the acceleration
scales as $\sqrt{\eta_{\rm coh}}$ rather than $\eta_{\rm coh}$---the
MOND square-root law halves the benefit of source enhancement.

\textbf{7. Deep-space advantage:}
Device performance improves as the background gravitational field
decreases.  The MOND crossover for the Solar System is at $\sim$7000~AU.
Beyond this, the MOND amplification factor for a local source grows as
$(a_0/a_{\rm bg})^{1/2}$.

\textbf{8. Bottom line for feasibility:}
The MOND crossover nonlinearity alone cannot close the $\sim 10^{40}$
gap between EM energy density and gravitational effects.  Quantum
coherence enhancement at $N$-scaling \emph{can} close the gap.  At
$\sqrt{N}$-scaling, the combination of MOND amplification + coherence
produces measurable effects ($\sim 10^{-11}$~m/s$^2$ with 1~MJ) but
requires very large stored energies for propulsion.

\textbf{9. The key experiment:}
Determining whether the EM$\to\psi$ coherence enhancement scales as
$\sqrt{N}$ or $N$ is the single most important measurement for the
propulsion application.  This can be tested with SRF cavities at
facilities like SQMS (Fermilab).

\end{tcolorbox}


%=============================================================================
\section{Appendix: Comparison of Interpolating Functions}\label{sec:appendix-mu}
%=============================================================================

For reference, we compare three interpolating functions used in the
MOND literature:

\begin{center}
\begin{tabular}{lccc}
\toprule
Function & $\mu(x)$ & $\mathcal{A}(1)$ & $d\mathcal{A}/dx|_{x=1}$ \\
\midrule
DFD canonical & $x/(1+x)$ & 2.00 & $-1.00$ \\
Standard & $x/\sqrt{1+x^2}$ & 1.414 & $-0.354$ \\
RAR (simple) & $1/(1+x^{-1})$ & same as DFD & same \\
Bekenstein & $x/(1+x^2)^{1/2}$ & 1.414 & $-0.354$ \\
\bottomrule
\end{tabular}
\end{center}

The DFD canonical function gives the strongest amplification at all
finite $x$, with the two families converging in both the Newtonian
($x \to \infty$) and deep-MOND ($x \to 0$) limits.

\begin{thebibliography}{10}
\bibitem{Alcock2025DFD} G.~T.~Alcock, ``Density Field Dynamics: A Complete
Unified Theory,'' (2025).
\end{thebibliography}

\end{document}
